\section{Introduction}
\label{sec:intro}

\subsection{Background}
Cables carry power and signals in almost every engineered system, from aircraft and ships to data centres, telecom networks and oil and geothermal wells. Over time they degrade through heating cycles, crushing, moisture ingress, insulation ageing and broken conductors. Each defect shows up electrically as a local change in the line's impedance. Finding such defects early, before they become shorts or outages, is the job of cable diagnostics.

Time-domain reflectometry (TDR) is the standard non-destructive way to do this \cite{pozar2011microwave,ulaby2015fundamentals}. A fast voltage step is launched into the cable, and the voltage $v(0,t)$ at the sending end is recorded. Wherever the impedance changes, part of the wave is reflected back. The arrival time of an echo gives the distance to the discontinuity, because the wave speed on a uniform line is constant. The height and sign of the echo describe the impedance change. This is the transmission-line theory of the EE~305 course applied directly: reflection coefficients, bounce diagrams and transients on terminated lines.

\subsection{Problem Statement}
Reading a TDR trace is easy for an ideal line and a single strong fault. Real measurements are harder for three reasons:
\begin{enumerate}
    \item \textbf{Loss and dispersion.} The conductor resistance grows with frequency through the skin effect ($R \propto \sqrt{f}$), and the dielectric adds a conductance $G \propto f$. Sharp edges are therefore rounded and attenuated as they travel, so distant echoes become small and slow and hard to time.
    \item \textbf{Overlapping reflections.} A fault and an imperfect load each reflect, and waves bounce between them. Their echoes can overlap, and a strong fault can hide the load entirely: a hard short or open ($\Gamma = \pm 1$) reflects all of the energy, so nothing reaches the load.
    \item \textbf{Noise.} Measurement noise hides weak discontinuities, such as a high-resistance shunt path or a small series resistance, whose reflection coefficient is only a few per cent.
\end{enumerate}
Machine learning has been applied to TDR since neural networks were first trained on simulated wiring-network reflectograms \cite{smail2011detection}. Convolutional networks have since been used to detect and characterise cable discontinuities \cite{scarpetta2021detection,novel2025wiring}, and image encodings of the trace \cite{goudjil2026wiring} and compressed-sensing TDR \cite{ieee2024cs_tdr} have been proposed for better noise robustness. Most of this work treats one task at a time, often on idealised lines. It rarely estimates the fault position and the load parameters ($R_L$, $L$, $C$) together.

\subsection{Objectives}
The goal of this project is to identify, from one noisy step-TDR trace of a lossy coaxial cable:
\begin{itemize}
    \item \textbf{Task T1 -- fault type} (7 classes): no fault (\texttt{none}), \texttt{short}, \texttt{open}, shunt resistance (\texttt{shunt\_R}), shunt capacitance (\texttt{shunt\_C}), series resistance (\texttt{series\_R}) and a short section of different impedance (\texttt{z0\_section}).
    \item \textbf{Task T2 -- fault location} $d$ along the cable, reported in metres.
    \item \textbf{Task T3 -- load type} (6 classes: \texttt{open}, \texttt{short}, \texttt{matched}, \texttt{R}, series \texttt{RL}, parallel \texttt{RC}), together with the load resistance $R_L$ and the time constant $\tau$, from which $L$ or $C$ follows.
\end{itemize}

\subsection{Approach and Deliverables}
The project has five parts, and Fig.~\ref{fig:pipeline} in Section~\ref{sec:methods} shows how they connect:
\begin{enumerate}
    \item a frequency-domain simulator of a lossy coaxial line, built on the telegrapher's equations and checked against analytical results;
    \item a labelled dataset of 85,000 traces, including five out-of-distribution test sets;
    \item three diagnostic methods: a conventional rule-based TDR detector, LightGBM \cite{ke2017lightgbm} on 94 engineered features, and a multi-task 1-D ResNet \cite{he2016deep};
    \item eight experiments on accuracy, noise, generalisation, detection limits, data efficiency, design choices and interpretability;
    \item a physics check on public laboratory measurements of a damaged 80~m coaxial cable \cite{nrel_wise_casing}.
\end{enumerate}
The complete code, data generator, trained models, an interactive demo app and this report are in the project repository. The full pipeline runs end to end with one command and is fully seeded, so every number in this report can be reproduced.
